\subsection{LAWS OF COMPOSITION}

DEFINITION 1. Let E be a set. A mapping f of E × E into E is called a law of composition on E. The value f(x, y) of f for an ordered pair (x, y) \ensuremath{\in} E × E is called the composition of x and y under this law. A set with a law of composition is called a magma.

The composition of x and y is usually denoted by writing x and y in a definite order and separating them by a characteristic symbol of the law in question (a symbol which it may be agreed to omit). Among the symbols most often used are + and ., the usual convection being to omit the latter if desired; with these symbols the composition of x and y is written respectively as \(x + y\) and x.y or xy. A law denoted by the symbol + is usually called addition (the composition \(x + y\) being called the sum of x and y) and we say that it is written additively; a law denoted by the symbol . is usually called multiplication (the composition x.y = xy being called the product of x and y) and we say that it is written multiplicatively. In the general arguments of paragraphs 1 to 3 of this chapter we shall generally use the symbols \(\top\) and \(\bot\) to denote arbitrary laws of composition.

By an abuse of language, a mapping of a subset of \(E \times E\) into E is sometimes called a law of composition not everywhere defined on E.

Examples. (1) The mappings \((\mathbf{X},\mathbf{Y})\mapsto \mathbf{X}\cup \mathbf{Y}\) and \((\mathbf{X},\mathbf{Y})\mapsto \mathbf{X}\cap \mathbf{Y}\) are laws of composition on the set of subsets of a set E.

\begin{enumerate}
\def\labelenumi{(\arabic{enumi})}
\setcounter{enumi}{1}
\item
  On the set \(\mathbf{N}\) of natural numbers, addition, multiplication and exponentiation are laws of composition (the compositions of \(x \in \mathbf{N}\) and \(y \in \mathbf{N}\) under these laws being denoted respectively by \(x + y\) , \(xy\) or \(x.y\) and \(x^y\) ) (Set Theory, III, § 3, no. 4).
\item
  Let E be a set; the mapping \((\mathbf{X}, \mathbf{Y}) \mapsto \mathbf{X} \circ \mathbf{Y}\) is a law of composition on the set of subsets of \(E \times E\) (Set Theory, II, § 3, no. 3, Definition 6); the mapping \((f, g) \mapsto f \circ g\) is a law of composition on the set of mappings of E into E (Set Theory, II, § 5, no. 2).
\item
  Let E be a lattice (Set Theory, III, § 1, no. 11); if \(\sup(x, y)\) denotes the least upper bound of the set \(\{x, y\}\) , the mapping \((x, y) \mapsto \sup(x, y)\) is a law of composition on E. Similarly for the greatest lower bound \(\inf(x, y)\) . Example 1 above is a particular case of this with \(\mathfrak{P}(E)\) ordered by inclusion.
\item
  Let \((\mathbf{E}_i)_{i\in I}\) be a family of magmas. Let \(\mathsf{T}_i\) denote the law of composition on \(\mathbf{E}_i\) . The mapping
\end{enumerate}

\[
((x _ {i}), (y _ {i})) \mapsto ((x _ {i} \top_ {i} y _ {i}))
\]

is a law of composition on the product \(E = \prod_{i \in I} E_i\) , called the product of the laws \(T_i\) . The set E with this law is called the product magma of the magmas \(E_i\) . In particular, if all the magmas \(E_i\) are equal to the same magma M, we obtain the magma of mappings of I from M.

Let \((x,y)\mapsto x\top y\) be a law of composition on a set E. Given any two subsets X, Y of E, X \(\top\) Y (provided this notation does not lead to confusion†) will denote the set of elements \(x\top y\) in E, where \(x\in X\) , \(y\in Y\) (in other words, the image of X \(\times\) Y under the mapping \((x,y)\mapsto x\top y\) ).

If \(a \in \mathbf{E}\) we usually write \(a \top \mathbf{Y}\) instead of \(\{a\} \top \mathbf{Y}\) and \(\mathbf{X} \top a\) instead of \(\mathbf{X} \top \{a\}\) . The mapping \((\mathbf{X}, \mathbf{Y}) \mapsto \mathbf{X} \top \mathbf{Y}\) is a law of composition on the set of subsets of \(\mathbf{E}\) .

DEFINITION 2. Let E be a magma and \(\top\) denote its law of composition. The law of composition \((x, y) \mapsto y \top x\) on E is called the opposite of the above. The set E with this law is called the opposite magma of E.

Let E and E' be two magmas; we shall denote their laws by the same symbol T. Conforming with the general definitions (Set Theory, IV, § 1, no. 5), a bijective mapping f of E onto E' such that

\[
f (x \top y) = f (x) \top f (y)\tag{1}
\]

for every ordered pair \((x, y) \in \mathbf{E} \times \mathbf{E}\) is called an isomorphism of \(\mathbf{E}\) onto \(\mathbf{E}'\) . \(\mathbf{E}\) and \(\mathbf{E}'\) are said to be isomorphic if there exists an isomorphism of \(\mathbf{E}\) onto \(\mathbf{E}'\) .

More generally:

DEFINITION 3. A mapping \(f\) of \(\mathbf{E}\) into \(\mathbf{E}'\) such that relation (1) holds for every ordered pair \((x, y) \in \mathbf{E} \times \mathbf{E}\) is called a homomorphism, or morphism, of \(\mathbf{E}\) into \(\mathbf{E}'\) ; if \(\mathbf{E} = \mathbf{E}'\) , \(f\) is called an endomorphism of \(\mathbf{E}\) .

The identity mapping of a magma E is a homomorphism, the composition of two homomorphism is a homomorphism.

For a mapping \(f\) of \(E\) into \(E'\) to be an isomorphism, it is necessary and sufficient that it be a bijective homomorphism and \(\overline{f}^{-1}\) is then an isomorphism of \(E'\) onto \(E\) .

